{
\begin{fileviewerbox}{Demonstration}
\begin{Verbatim}[breaklines=true]
Here is a demonstration of how to correctly accomplish this task.
It is included to show you how to correctly use the interface.
You do not need to follow exactly what is done in the demonstration.
--- DEMONSTRATION ---
{{INSTANCE PROMPT}}

The code has been updated to use the `round` function, which should fix the rounding issue. Before submitting the changes, it would be prudent to run the reproduce.py code again to ensure that this change has the desired effect. I see that "Current directory" is still `/marshmallow-code__marshmallow`, so we can just run `reproduce.py` directly.
```
python reproduce.py
```
345

(Open file: /marshmallow-code__marshmallow/src/marshmallow/fields.py)
(Current directory: /marshmallow-code__marshmallow)
bash-$
The output has changed from 344 to 345, which suggests that the rounding issue has been fixed. Let's remove the reproduce.py file since it is no longer needed.
```
rm reproduce.py
```
Your command ran successfully and did not produce any output.
(Open file: /marshmallow-code__marshmallow/src/marshmallow/fields.py)
(Current directory: /marshmallow-code__marshmallow)
bash-$
rm doesn't have any output when it deletes successfully, so that must have worked. Now that we have fixed the issue, we should submit our changes to the code base. Let's submit the changes using the `submit` command.
```
submit
```
--- END OF DEMONSTRATION ---
\end{Verbatim}
\end{fileviewerbox}
\captionsetup{type=figure}
\captionof{figure}{
A simplified demonstration template showing how demonstrations are provided to the model as a single message.
Here we show only the final $3$ turns in the demonstration for brevity.
}
\label{fig.prompts.demonstration}
}